\documentclass{article}
\usepackage{graphicx}
\usepackage{hyperref}
\usepackage{booktabs}

\title{An Analysis on the Long Tail Impacts of Subprime Loans in the Early 2000s}
\author{Matthew Robson}
\date{July 2026}

\begin{document}

\maketitle

\section{Introduction}

\subsection{Subprime Loans}

Subprime Loans have been the focus of many papers on the 2008 financial crisis but are often disagreed on when it comes to specific definitions \cite{mian2009consequences}. Generally, subprime loans are considered loans given to borrowers of high risk. This generally means that these loans will have higher interest rates or stricter requirements such as in the down payment or penalties. For the purposes of this paper, we will consider any loan with a rate spread greater than or equal to 3.00\% as subprime. This has some drawbacks in cases where the rate spread may not indicate some other requirements of the contract or the loan or may be measured inaccurately.  

Some argue that origins of subprime loans are in adjustable rate mortgages---mortgages where the interest rate for the borrower may change based on external economic circumstances \cite{duke_predatory_lending}. Adjustable rate mortgages often lead to a predatory structure due to the increased rates over time. This coincides with subprime loans because both generally have harsher terms and so are taken out by people that have poor credit. The prevalence of subprime loans began in the early 2000s and held up to the financial crisis of 2008, after which, lending was significantly decreased. This increase in subprime loans during the early 2000s is believed to have driven an increase in construction of single family home in comparison to multi-unit housing. This is because many people that previously did not have access to taking out a loan for the construction of their own house now were given the opportunity as more loan agencies were offering subprime loans to a much wider set of people.

Subprime loans like these were then traded on the public market in many different forms such as collateralized debt obligations (CDOs) and mortgage-backed securities (MBS). A CDO is a financial device that consists of many mortgages and spreads the risk into "tranches" with differing return rates and risk rates. Generally the rating scheme for CDOs is considered highly misleading because while the senior tranches might be more secure than the mezzanine or equity, the senior tranche is still backed by loans that may default leading to loss of profits. MBS are very similar but only directed at housing loans. 

\subsection{The 2008 Financial Crisis}

The 2008 financial crisis is often attributed to the over reliance on subprime loans and over-leveraging of the early 2000s. Because many firms in the 

\subsection{The Current Housing Crisis}

The current housing market is full of difficulty for new and low income buyers. Average house prices have been rising relative to average income for the past few years, leading to a high cost burden on those wanting to purchase a house. In some cases nearly 70\% of a household's income is used directly on paying rent \cite{nlihc_problem}. In addition to this, housing has become overcrowded such that many people are forced to live in cramped conditions \cite{naeh_state_of_homelessness}. There have been many proposed causes for this crisis such as the over indebtedness of people today, rapid urbanization outpacing housing development. 

\subsection{The Impacts of Subprime Loans}

\section{Methods}

\subsection{Data}

We collected loan data for this study from the HMDA. From our collected data, we used Loan data from the HMDA from 2004 to 2007. Data from 2000 to 2006 was sourced from OPENICPSR \cite{openicpsr_hmda_2000_2006}, and data from 2007 - 2010 was sourced from the Consumer Financial Protection Bureau website \cite{cfpb_hmda_2007_2010}. This data was used to generate the percent subprime loans for each census tract over this time period. Specifically, we focused on the rate spread variable to classify loans as subprime or prime. We used the ACS 5 year estimates \cite{census_acs_2016_2020} to obtain the data for our dependent variables and control variables. We used data from the Housing, Socioeconomic, and Demographic data sets for this analysis. All of the ACS data came from the 2016 - 2020 time frame. Some of the models included use spatial data for the shapes of the census tracts, and this data came from IPUMS NHGIS \cite{ipums_nhgis}. 



\subsection{Dependent Variables}

The dependent variable considered for this study was housing insecurity. While housing insecurity is a complex topic with many factors that make it up, we used cost burden and overcrowding as a proxy for housing insecurity. In this study, we consider someone to be "cost burdened" by their housing if greater than 35\% of their income has gone to housing in the past 12 months. We consider the overcrowding rate of a census tract to be calculated with the following formula (Housing Units With Greater than 1 Person per Room) / (Number of Housing Units).   

\subsection{Lending Variable}

The variable used for lending was rate spread. Rate spread is defined as the loan interest rate minus the average loan interest rate. We classified any value above 3\% as subprime lending and any value below 3\% as prime lending. 

\subsection{Control Variables}

\section{Results}

\subsection{Summary Statistics}

\subsection{Regression Table 1}

\begin{table}[htbp]
\centering
\caption{Regression Results: Overcrowding Rate}
\label{tab:reg_overcrowding}
\small
\begin{tabular}{lcccc}
\toprule
 & \multicolumn{4}{c}{\textit{Dependent Variable:} Overcrowding Rate} \\
\cmidrule(lr){2-5}
 & (1) Baseline & (2) Economic & (3) Demographic & (4) All Variables \\
\midrule
Subprime Loan (\%)\quad & 0.1562*** & 0.1029*** & 0.0062     & 0.0024 \\
                        & (0.0048)  & (0.0065)  & (0.0074)   & (0.0073) \\[0.5em]
Unemployment Rate (\%)  &           & -0.0002** &            & -0.0008*** \\
                        &           & (0.0001)  &            & (0.0001) \\[0.5em]
Poverty Rate (\%)       &           & 0.0009*** &            & 0.0002*** \\
                        &           & (0.0000)  &            & (0.0000) \\[0.5em]
Median Household Income &           & 0.0001*** &            & 0.0001*** \\
                        &           & (0.0000)  &            & (0.0000) \\[0.5em]
White (\%)              &           &           & -0.0013*** & -0.0013*** \\
                        &           &           & (0.0000)   & (0.0000) \\[0.5em]
Black (\%)              &           &           & -0.0012*** & -0.0011*** \\
                        &           &           & (0.0000)   & (0.0000) \\[0.5em]
Hispanic (\%)           &           &           & -0.0002*** & -0.0002*** \\
                        &           &           & (0.0000)   & (0.0000) \\[0.5em]
Undergraduate (\%)      &           &           & -0.0004*** & -0.0005*** \\
                        &           &           & (0.0000)   & (0.0000) \\[0.5em]
Constant                & 0.0138*** & 0.0039*** & 0.1438***  & 0.1425*** \\
                        & (0.0005)  & (0.0012)  & (0.0028)   & (0.0030) \\
\midrule
Observations   & 44,611 & 44,489 & 44,611 & 44,489 \\
$R^2$          & 0.0270 & 0.0581 & 0.3968 & 0.4044 \\
Adjusted $R^2$ & 0.0270 & 0.0580 & 0.3967 & 0.4043 \\
\bottomrule
\multicolumn{5}{l}{\footnotesize Standard errors in parentheses.} \\
\multicolumn{5}{l}{\footnotesize * $p < 0.1$, ** $p < 0.05$, *** $p < 0.01$}
\end{tabular}
\end{table}

\subsection{Regression Table 2}

\begin{table}[htbp]
\centering
\caption{Regression Results: Cost Burden}
\label{tab:reg_cost_burden}
\small
\begin{tabular}{lcccc}
\toprule
 & \multicolumn{4}{c}{\textit{Dependent Variable:} Cost Burden Rate} \\
\cmidrule(lr){2-5}
 & (1) Baseline & (2) Economic & (3) Demographic & (4) All Variables \\
\midrule
Subprime Loan (\%)      & -36.8989*** & -12.3028*** & -19.8165*** & -25.1691*** \\
                        & (1.0786)    & (1.4482)    & (1.7683)    & (1.6327) \\[0.5em]
Unemployment Rate (\%)  &             & 0.2010***   &             & -0.0043 \\
                        &             & (0.0158)    &             & (0.0156) \\[0.5em]
Poverty Rate (\%)       &             & -0.0059     &             & -0.2112*** \\
                        &             & (0.0075)    &             & (0.0082) \\[0.5em]
Median Household Income &             & 0.0786***   &             & 0.0266*** \\
                        &             & (0.0021)    &             & (0.0022) \\[0.5em]
White (\%)              &             &             & -0.2518***  & -0.2754*** \\
                        &             &             & (0.0065)    & (0.0066) \\[0.5em]
Black (\%)              &             &             & -0.1558***  & -0.1247*** \\
                        &             &             & (0.0071)    & (0.0064) \\[0.5em]
Hispanic (\%)           &             &             & -0.0795***  & -0.0678*** \\
                        &             &             & (0.0078)    & (0.0067) \\[0.5em]
Undergraduate (\%)      &             &             & 0.1269***   & 0.0525*** \\
                        &             &             & (0.0036)    & (0.0039) \\[0.5em]
Constant                & 20.2390***  & 10.9132***  & 32.2269***  & 37.7072*** \\
                        & (0.1329)    & (0.3008)    & (0.7134)    & (0.6408) \\
\midrule
Observations   & 44,716 & 44,489 & 44,716 & 44,489 \\
$R^2$          & 0.0291 & 0.0722 & 0.2389 & 0.2981 \\
Adjusted $R^2$ & 0.0291 & 0.0721 & 0.2389 & 0.2980 \\
\bottomrule
\multicolumn{5}{l}{\footnotesize Standard errors in parentheses.} \\
\multicolumn{5}{l}{\footnotesize * $p < 0.1$, ** $p < 0.05$, *** $p < 0.01$}
\end{tabular}
\end{table}

\section{Discussion}

\bibliographystyle{plain}
\bibliography{references}

\end{document}


